\section{Model}

\begin{frame}{Your Solution}
  \begin{center}
    \resizebox{\textwidth}{!}{%
    \begin{tikzpicture}[node distance=1.5cm, every node/.style={font=\small}]
      \node[draw, rounded corners, minimum width=2.8cm, minimum height=0.9cm] (input) {Telemetry window};
      \node[draw, rounded corners, right=of input, minimum width=2.5cm, minimum height=0.9cm] (enc) {Encoder};
      \node[draw, rounded corners, right=of enc, minimum width=3cm, minimum height=0.9cm] (loop) {Looped latent core};
      \node[draw, rounded corners, right=of loop, minimum width=2.8cm, minimum height=0.9cm] (head) {Prediction heads};
      \draw[-{Latex}] (input) -- (enc);
      \draw[-{Latex}] (enc) -- (loop);
      \draw[-{Latex}] (loop) -- (head);
      \draw[-{Latex}, bend left=35] (loop.north) to node[above]{exit gate} (loop.north east);
    \end{tikzpicture}
    }%
  \end{center}
  \vspace{0.3cm}
  \begin{itemize}
    \item First, we turn the sensor window into a compact representation.
    \item Then the looped model can refine its answer step by step.
    \item If it becomes confident early, it stops early.
    \item Final outputs: machine health, maintenance action, and likely subsystem.
  \end{itemize}
\end{frame}

\begin{frame}{Why This Is Different From the LLM Baseline}
  \begin{itemize}
    \item The looped model is built for numeric sensor streams.
    \item The LLM baseline has to read the same data after it is converted into text.
    \item That makes this a fairer question than ``looped model vs generic chatbot''.
    \item We tested whether the best task-specific version of each approach works better.
  \end{itemize}
\end{frame}
